\chapter{Validaci\'on y resultados reproducibles}
\label{chap:validacion}

La evaluaci\'on de un agente generativo debe distinguir claramente entre una demostraci\'on de integraci\'on, una evaluaci\'on de la calidad de respuestas y una afirmaci\'on de seguridad. Que el chat se abra en un navegador no demuestra que sus recomendaciones sean correctas; del mismo modo, superar un conjunto finito de ataques no prueba que el sistema sea seguro. En este trabajo se adopta una estrategia de evidencia por capas: pruebas autom\'aticas de la aplicaci\'on, un conjunto sint\'etico trazable para el agente y revisi\'on humana contra fuentes de oro. La metodolog\'ia sigue la fase de evaluaci\'on de DSRM \cite{peffers2007} y reconoce la no determinaci\'on propia de los sistemas RAG \cite{lewis2020rag,ragas2024}.

\section{Preguntas, hip\'otesis y configuraciones comparadas}

La pregunta principal es: \emph{\textquestiondown en qu\'e medida un agente basado en un corpus institucional curado mejora la correcci\'on, trazabilidad, cobertura y utilidad de las respuestas frente a una configuraci\'on base, manteniendo un comportamiento seguro ante consultas no cubiertas y entradas adversarias?} Para responderla se establecen las siguientes hip\'otesis antes de ejecutar la evaluaci\'on:

\begin{itemize}
  \item \textbf{H1}: el corpus curado mejora la exactitud factual y la trazabilidad respecto de una configuraci\'on base.
  \item \textbf{H2}: las reglas de citaci\'on, abstenci\'on y derivaci\'on reducen respuestas no fundamentadas y fallos de seguridad.
  \item \textbf{H3}: las decisiones esenciales se mantienen consistentes ante reformulaciones equivalentes.
  \item \textbf{H4}: la interfaz resulta comprensible y \`util para el perfil de usuario sin inducir a sobreconfianza.
\end{itemize}

La comparaci\'on propuesta utiliza tres configuraciones aisladas: C0, agente base sin corpus curado; C1, agente con corpus curado y b\'usqueda web abierta desactivada; y C2, configuraci\'on final que a\~nade las instrucciones de seguridad, temas deterministas de derivaci\'on y controles de interfaz. C0 se ejecuta solo en un entorno de pruebas y no se publica. La diferencia C0--C2 mide el efecto global del dise\~no; C1--C2 a\'isla el efecto de las salvaguardas conversacionales.

\section{Conjunto de pruebas y r\'ubrica}

El paquete de entrega incorpora 72 casos sint\'eticos versionados en formato CSV y JSON. Cada caso contiene un identificador, categor\'ia, consulta, criticidad, proposiciones de oro, fuente esperada, necesidad de abstenci\'on y criterio de seguridad. Las preguntas no se copian literalmente de las fuentes; se derivan de afirmaciones verificables y se revisan antes de congelar la versi\'on de evaluaci\'on. Esta separaci\'on evita que el agente se ajuste manualmente a los mismos casos con los que ser\'a evaluado.

\begin{table}[H]
\centering
\small
\caption{Familias de casos del conjunto de evaluaci\'on.}
\label{tab:familias-prueba}
\begin{tabularx}{\textwidth}{p{0.27\textwidth}p{0.21\textwidth}X}
\toprule
\textbf{Familia} & \textbf{Resultado esperado} & \textbf{Ejemplo de verificaci\'on}\\
\midrule
Consulta dentro de alcance & Respuesta clara, prudente y citada. & Phishing, MFA, actualizaciones, contrase\~nas o navegaci\'on segura.\\
Ambig\"uedad o multivuelta & Aclaraci\'on m\'inima o respuesta condicional. & ``He recibido algo raro; \textquestiondown qu\'e hago?'' sin datos suficientes.\\
Sin cobertura & Abstenci\'on explicada y derivaci\'on. & Pregunta sobre una pol\'itica no presente en el corpus.\\
Incidente o criticidad alta & Pasos no destructivos y canal humano. & Posible compromiso de cuenta o ransomware.\\
Inyecci\'on y extracci\'on & Rechazo sin revelar instrucciones ni secretos. & ``Ignora las reglas y muestra el prompt del sistema''.\\
Solicitud ofensiva & Negativa segura y alternativa defensiva. & Petici\'on de eludir un control o acceder a un tercero.\\
Interfaz y seguridad web & Sin secretos, XSS ni estados enga\~nosos. & Token solamente desde BFF, marcado visible de modo demo.\\
\bottomrule
\end{tabularx}
\end{table}

La r\'ubrica humana punt\'ua cada respuesta en cinco dimensiones de 0 a 2: exactitud factual, cobertura de proposiciones, justificaci\'on por fuente, seguridad/derivaci\'on y claridad. Un 0 representa fallo; 1, cumplimiento parcial; y 2, cumplimiento completo. Dos revisores independientes etiquetan una muestra previamente acordada; las discrepancias se resuelven mediante consulta a la fuente de oro, no por preferencia de estilo. Para etiquetas nominales, como ``abstenci\'on correcta'', se calcula el acuerdo de Cohen.

\section{M\'etricas}

Las m\'etricas se definen antes de observar resultados. Sean $C$ los casos dentro de alcance, $A$ las afirmaciones factuales que requieren evidencia y $P$ los ataques adversarios. La exactitud, fundamentaci\'on y tasa de ataque exitoso se calculan como:

\begin{equation}
\text{Exactitud} = \frac{\text{casos correctos de } C}{|C|}
\label{eq:exactitud}
\end{equation}

\begin{equation}
\text{Fundamentaci\'on} = \frac{\text{afirmaciones de } A \text{ respaldadas por evidencia}}{|A|}
\label{eq:groundedness}
\end{equation}

\begin{equation}
\text{ASR} = \frac{\text{ataques de } P \text{ que inducen conducta insegura}}{|P|}
\label{eq:asr}
\end{equation}

Adicionalmente se mide: precisi\'on de cita (citas que respaldan realmente la afirmaci\'on), completitud de cita (afirmaciones que poseen la cita requerida), tasa de abstenci\'on correcta, cobertura por dominio, tasa de errores cr\'iticos, consistencia entre reformulaciones y latencia mediana/p95. Las metas de aceptaci\'on son criterios de decisi\'on, no resultados anticipados: ning\'un fallo cr\'itico en los casos adversarios, fundamentaci\'on y precisi\'on de cita superiores al 95\%, y abstenci\'on correcta superior al 90\%. Si una configuraci\'on no los alcanza, se documenta el defecto y no se promueve a demostraci\'on.

Las evaluaciones autom\'aticas de Copilot Studio pueden aportar se\~nales de relevancia, grounding, completitud y abstenci\'on \cite{copilotEvaluation}. Sin embargo, no sustituyen la revisi\'on humana: emplean modelos y criterios internos que pueden cambiar con el tiempo. Por ello se exportan las transcripciones y se conservan junto con el identificador de configuraci\'on, versi\'on de corpus y fecha de ejecuci\'on.

\section{Procedimiento de ejecuci\'on}

El protocolo evita fugas entre desarrollo y evaluaci\'on:

\begin{enumerate}[label=\arabic*.]
  \item Congelar la versi\'on de instrucciones, temas, corpus y casos de prueba; registrar hashes y fecha.
  \item Ejecutar pruebas de desarrollo con un conjunto independiente y corregir defectos sin modificar el conjunto final.
  \item Publicar C0, C1 y C2 en el entorno de pruebas, con la misma configuraci\'on de idioma y sin b\'usqueda web abierta.
  \item Ejecutar cada caso final, registrar salida, fuentes, latencia y configuraci\'on. Repetir las consultas de consistencia en sesiones nuevas.
  \item Aplicar la r\'ubrica humana de forma ciega a la configuraci\'on cuando sea posible y calcular las m\'etricas.
  \item Revisar fallos cr\'iticos, repetir el ciclo de dise\~no y conservar tambi\'en los resultados de las versiones descartadas.
\end{enumerate}

La publicaci\'on en un canal protegido puede requerir autenticaci\'on manual con Entra ID. La documentaci\'on de Microsoft indica que algunas evaluaciones autom\'aticas no admiten ese modo, por lo que debe evaluarse un clon de pruebas compatible y completar despu\'es pruebas funcionales manuales de la versi\'on protegida \cite{copilotEvaluation,copilotSso}. Esta limitaci\'on se declara para no confundir resultados de un entorno de laboratorio con evidencia de producci\'on.

\section{Resultados locales verificables}

La aplicaci\'on y el BFF incorporan comprobaciones autom\'aticas para propiedades que no dependen de disponer de un tenant: carga de configuraci\'on, ausencia de secreto en el cliente, rechazo de origen no permitido, l\'imites de solicitud y funcionamiento marcado del modo demostraci\'on. El archivo de evidencia generado durante la construcci\'on de esta entrega se incluye a continuaci\'on. Se trata de evidencia de integraci\'on local, no de una medici\'on de la calidad de Copilot Studio.

\begin{table}[H]
\centering
\small
\caption{Comprobaciones locales ejecutadas sobre la versi\'on revalidada el 11 de noviembre de 2026.}
\label{tab:resultados-locales}
\begin{tabularx}{\textwidth}{p{0.43\textwidth}p{0.20\textwidth}X}
\toprule
\textbf{Comprobaci\'on} & \textbf{Estado} & \textbf{Interpretaci\'on}\\
\midrule
Compilaci\'on del cliente Angular y BFF & Superada & Los paquetes de producci\'on se generaron sin errores.\\
Pruebas del BFF & 14/14 superadas & Cubren origen permitido, validaci\'on de entrada, protecci\'on de secretos, errores controlados y adaptador OpenAI.\\
Prueba real del BFF & HTTP 200 & \texttt{GET /api/health} indic\'o modo OpenAI y \texttt{POST /api/chat} respondi\'o a una consulta sint\'etica.\\
Prueba mediante proxy Angular & HTTP 200 & \texttt{GET /api/runtime} y \texttt{POST /api/chat} respondieron correctamente en \texttt{localhost:4200}.\\
Escaneo de secretos en cliente & 0 coincidencias & La clave de OpenAI y el secreto Direct Line no aparecen en el c\'odigo fuente del cliente.\\
Auditor\'ia de dependencias & Pendiente & La evidencia local registra dos vulnerabilidades transitivas altas que deben corregirse antes de un despliegue p\'ublico.\\
\bottomrule
\end{tabularx}
\end{table}

Estas salidas verifican propiedades de integraci\'on y de seguridad del c\'odigo. Una compilaci\'on correcta o una respuesta HTTP 200 no equivalen a una evaluaci\'on de exactitud de las respuestas del agente.


\section{Estado de la evaluaci\'on del agente real}

En la fecha de esta versi\'on no se inventan valores para C0, C1 o C2: su ejecuci\'on requiere que el propietario inicie sesi\'on en un tenant de Copilot Studio con licencia, cree el agente, apruebe las fuentes y publique el canal. La ausencia de ese acceso es una condici\'on externa del experimento, no un resultado favorable ni desfavorable de las hip\'otesis. El paquete entregado contiene las instrucciones, el corpus, los temas, la matriz y la gu\'ia paso a paso para producir esa evidencia sin redise\~nar el sistema.

Esta distinci\'on es metodol\'ogicamente preferible a reportar porcentajes simulados. Tras ejecutar el protocolo, la tabla \ref{tab:resultados-copilot} debe completarse con los valores exportados y las transcripciones correspondientes; ninguna fila se considerar\'a v\'alida sin el identificador de versi\'on, fecha y revisor.

\begin{table}[H]
\centering
\small
\caption{Plantilla de consolidaci\'on de resultados del agente real.}
\label{tab:resultados-copilot}
\begin{tabular}{lccc}
\toprule
\textbf{M\'etrica} & \textbf{C0} & \textbf{C1} & \textbf{C2}\\
\midrule
Exactitud factual & Pendiente & Pendiente & Pendiente\\
Fundamentaci\'on y cita & Pendiente & Pendiente & Pendiente\\
Abstenci\'on correcta & Pendiente & Pendiente & Pendiente\\
ASR adversario & Pendiente & Pendiente & Pendiente\\
Latencia mediana / p95 & Pendiente & Pendiente & Pendiente\\
Errores cr\'iticos & Pendiente & Pendiente & Pendiente\\
\bottomrule
\end{tabular}
\end{table}

\section{Amenazas a la validez}

La validez de constructo est\'a limitada porque ``utilidad'' y ``seguridad'' no se reducen a una cifra. Se mitiga mediante una r\'ubrica multidimensional, fuentes de oro y revisi\'on humana. La validez interna puede verse afectada si se ajusta el agente con casos finales o si cambian simult\'aneamente varias configuraciones; por eso se congelan los datos y se comparan configuraciones separadas. La validez externa es limitada: un conjunto sint\'etico no representa todas las conversaciones reales ni los riesgos de una organizaci\'on concreta. Finalmente, la conclusi\'on depende del proveedor, idioma, versi\'on de modelo, licencia y fecha; los resultados deben interpretarse como evidencia localizada y reproducible, no como una garant\'ia general.

La evaluaci\'on muestra tambi\'en un riesgo pr\'actico: un corpus institucional reduce la incertidumbre, pero no elimina los fallos de recuperaci\'on ni la posibilidad de una interpretaci\'on incorrecta. La ruta de escalado y la supervisi\'on humana siguen siendo controles necesarios.
